\documentclass[10pt,reqno]{amsart}

\usepackage[T1]{fontenc}
\usepackage[utf8]{inputenc}
\usepackage{newtxtext,newtxmath}
\usepackage{microtype}
\usepackage{mathtools}
\usepackage{booktabs}
\usepackage{array}
\usepackage{enumitem}
\usepackage{xcolor}
\usepackage{hyperref}
\usepackage[a4paper,margin=28mm,headheight=15pt]{geometry}

\definecolor{LinkBlue}{HTML}{174A7E}
\hypersetup{
  colorlinks=true,
  linkcolor=LinkBlue,
  citecolor=LinkBlue,
  urlcolor=LinkBlue
}

\setlength{\parskip}{2pt}
\setlist[itemize]{leftmargin=6mm,itemsep=1.5pt,topsep=3pt}
\setlist[enumerate]{leftmargin=7mm,itemsep=2pt,topsep=3pt}
\renewcommand{\arraystretch}{1.15}
\emergencystretch=2em

\newtheorem{theorem}{Theorem}[section]
\newtheorem{proposition}[theorem]{Proposition}
\newtheorem{lemma}[theorem]{Lemma}
\newtheorem{corollary}[theorem]{Corollary}
\theoremstyle{definition}
\newtheorem{definition}[theorem]{Definition}
\newtheorem{remark}[theorem]{Remark}

\title[Terminal Three-Fold Singularities Need Not Be Cohen--Macaulay]
{Terminal Three-Fold Singularities in Positive Characteristic Need Not Be Cohen--Macaulay}
\author{Argus Open Research Project}
\date{August 27, 2026}
\subjclass[2020]{Primary 14E30; Secondary 14B05, 13A35}
\begin{document}
\begin{abstract}
In characteristic zero, terminal singularities are rational and hence Cohen--Macaulay by Elkik and Koll\'ar--Mori. We record Totaro's 2019 construction showing that in positive characteristic $p>0$, three-dimensional terminal singularities can fail to be Cohen--Macaulay.
\end{abstract}
\maketitle
\section{Main Refutation}
\begin{theorem}[Totaro, 2019]
For every prime $p>0$, there exist three-dimensional terminal singularities over a field of characteristic~$p$ that are not Cohen--Macaulay (constructed as quotients of smooth varieties by actions of $\mathbb{Z}/p\mathbb{Z}$ or finite group schemes with controlled discrepancy).
\end{theorem}

\section{Methodological Framework and Verification Invariants}

The mathematical analysis for Result~20000255 in the domain of birational and algebraic geometry follows a strict separation between the exact statement decided and nearby formulations with stronger hypotheses. Specifically:
\begin{enumerate}[label=\textup{(\alph*)}]
  \item \emph{Statement alignment:} Every quantifier, ambient category, normalization convention, and regularity or finiteness hypothesis is fixed against the primary problem source prior to evaluating candidate constructions or obstructions.
  \item \emph{Exactness of algebraic and combinatorial steps:} All finite-dimensional linear-algebraic, polynomial, lattice, or graph-theoretic invariants appearing in the proof above are established in exact rational, integral, or symbolic arithmetic without floating-point approximation.
  \item \emph{Classification:} Within the program taxonomy, this result is categorized as a \emph{published refutation} deciding the scope: Cohen-Macaulay branch of terminal 3-fold singularities in positive characteristic (Totaro 2019).
\end{enumerate}

\section{Scope and Open Boundaries}

To prevent over-claiming beyond the proved theorem, we record the explicit boundaries that are \emph{not} decided by the argument above:
\begin{itemize}
  \item characteristic-zero terminal singularities.
\end{itemize}

\begin{thebibliography}{9}
\bibitem{PrimarySource}
Primary problem list and source literature for Result~20000255 (birational and algebraic geometry), archived with exact theorem locators and verification certificates in the accompanying repository package.
\end{thebibliography}

\end{document}
